\paragraph{LoRA and its initialisation.} A LoRA adapter replaces a weight
update by $\Delta W = BA$ with $B \in \R^{m\times r}$, $A \in \R^{r\times d}$
and $r \ll \min(m, d)$~\citep{hu2022lora}. In the reference implementation and
in PEFT, $B$ starts at zero and $A$ is drawn from a Kaiming-uniform
distribution using the global random generator. The two factors play different
roles: $A$ mostly selects input features and $B$ produces the output, and a
random untrained $A$ performs almost as well as a trained
one~\citep{zhu2024asymmetry}. Several designs share $A$ on purpose.
HydraLoRA~\citep{tian2024hydralora} uses one $A$ for several $B$ heads, and
federated variants freeze $A$ (FFA-LoRA~\citep{sun2024ffalora}) or aggregate
only $A$ (FedSA-LoRA~\citep{guo2025fedsalora}) to make aggregation exact.
Other initialisations reach a shared $A$ without a shared seed:
PiSSA~\citep{meng2024pissa}, OLoRA~\citep{buyukakyuz2024olora} and
MiLoRA~\citep{wang2025milora} compute $A$ from the base weight, so every task
fine-tuned from one base with the same modules and rank starts from the same
$A$. CorDA~\citep{yang2024corda} orients its decomposition by task data, and
LoRA-FA~\citep{zhang2023lorafa} freezes a random $A$, so the first starts each
task differently and the second is shared only when the seed is. This
literature has not examined what such cohorts do to a merge.

\paragraph{Merging methods, by what they compute.} It helps to sort merging
methods by the operation they apply to the adapters. \emph{Averaging and
sparsifying} methods form a weighted combination of the task updates: task
arithmetic~\citep{ilharco2023task}, TIES~\citep{yadav2023ties},
DARE~\citep{yu2024dare}, DELLA~\citep{della2024} and the quantised
TVQ~\citep{kim2025tvq}. \emph{Subspace} methods act on the singular structure
of the updates: KnOTS~\citep{stoica2025knots} aligns them in a shared basis
and TSV-M~\citep{gargiulo2025tsv} whitens their stacked singular vectors.
\emph{Solvers} find the merged weight as the solution of a linear system.
RegMean~\citep{jin2023regmean} solves a least-squares problem in per-task
input-activation Gram matrices, MaTS~\citep{tam2024mats} solves a system in a
Fisher surrogate by conjugate gradient, and RegMean++~\citep{regmeanpp2025}
and LoRM~\citep{lorm2024} extend the closed form. Fisher
merging~\citep{matena2022fisher} weights each parameter by a diagonal Fisher
estimate. For this paper, the decisive property of a solver is where its
matrix comes from: from data, as in RegMean and MaTS, or from the adapters'
own factors, as in the data-free form of RegMean that replaces each activation
Gram by $A_t^\top A_t$.

\paragraph{Analyses of merging.} Weight
disentanglement~\citep{ortizjimenez2023disentanglement} identifies the
low-rank curvature structure of LoRA fine-tunes that the quantities below
adopt. Cao et al.~\citep{cao2026collapse} report that parameter-space conflict metrics
predict little about when merging fails, while representational
incompatibility between tasks predicts a great deal, and they prove a limit
for convex merges. Our results agree where the two overlap: we also find that
parameter-space geometry does not affect averaging-type merges
(\S\ref{sec:merging}). We add that one parameter-space quantity does predict
the behaviour of solvers built from the adapters, which are not convex
combinations of the fine-tuned models. A large empirical
study~\citep{systematic2025merging} finds that only task arithmetic reliably
produces constructive interference across four model families; that result
concerns which methods work on average, and ours concerns a property of the
cohort that decides one class of them.

\paragraph{Notation.} A cohort has $T$ adapters $\{(A_t, B_t)\}$ of rank $r$ on
a common base. The input subspace of adapter $t$ in a layer is the row space
$V_t = \range(A_t^\top) \subset \R^d$, and $P_t$ is the orthogonal projector
onto it. We measure overlap between task subspaces by the cosines of the
principal angles between each pair $V_i, V_j$, and report their median over
pairs and over a sample of layers. Following the weight-disentanglement
model, the operator a solver must invert is
\begin{equation}\label{eq:hbar}
  \bar H = \frac1T \sum_{t=1}^T P_t,
\end{equation}
and the exact interpolator, the minimum-norm update that reproduces every task
update on its own subspace, is $\tauH = \bar D\,\bar H^{+}$ with
$\bar D = \tfrac1T\sum_t \Delta W_t$ and $(\cdot)^+$ the pseudoinverse. We
report the condition number $\kappa(\bar H)$, the ratio of its largest to its
smallest nonzero eigenvalue, and the \emph{amplification}
$\|\tauH\| / \|\bar D\|$, how much larger the exact interpolator is than the
mean update. Both are computed from the adapter weights alone.
